\section*{Bab \ref{ch:b1:central-forces} --- Gaya Sentral: Planet dan Satelit}

\begin{solution}{exo:b1:central-forces:1}
$r = \qty{6770}{km}$: $v = \sqrt{GM/r} = \qty{7.7}{km/s}$; $T = 2\pi r/v
= \qty{5.5e3}{s} = \qty{92}{min}$; $1440/92 \approx 16$ kali matahari terbit tiap hari.
\end{solution}

\begin{solution}{exo:b1:central-forces:2}
$r = (GMT^2/4\pi^2)^{1/3} = (3.99\times10^{14} \times 7.42\times10^9/39.5)^{1/3}
= \qty{4.22e7}{m}$: ketinggiannya \qty{35800}{km}; $v = 2\pi r/T = \qty{3.1}{km/s}$.
\end{solution}

\begin{solution}{exo:b1:central-forces:3}
$v_e = \sqrt{2GM/R}$: Bumi \qty{11.2}{km/s}; Bulan $\sqrt{2 \times 4.90\times
10^{12}/1.74\times10^6} = \qty{2.4}{km/s}$. Molekul gas yang berkelajuan
beberapa ratus meter tiap sekon mempunyai ekor distribusi kelajuan di
atas \qty{2.4}{km/s}; sepanjang zaman geologi, gas Bulan bocor pergi.
\end{solution}

\begin{solution}{exo:b1:central-forces:4}
$T = a^{3/2}$ (tahun, AU): Mars \qty{1.88}{yr}; Jupiter \qty{11.9}{yr};
\qty{30}{AU}: \qty{164}{yr}.
\end{solution}

\begin{solution}{exo:b1:central-forces:5}
$\Delta E = \tfrac12 mv^2 + GMm(1/R - 1/r) = 2.94\times10^{10} + 3.99\times
10^{17} \times (1.570 - 1.477)\times10^{-7} = 2.94\times10^{10} + 0.37\times
10^{10} = \qty{3.3e10}{J}$ --- \qty{33}{MJ/kg}, yaitu energi \qty{0.7}{kg}
bensin tiap kilogram yang ditempatkan di orbit (roket memerlukan jauh
lebih banyak, karena harus mengangkat bahan bakarnya sendiri).
\end{solution}

\begin{solution}{exo:b1:central-forces:6}
$a = (0.586 + 35.1)/2 = \qty{17.8}{AU}$; $e = (35.1 - 0.586)/35.7 = 0.967$;
$T = 17.8^{3/2} = \qty{75}{yr}$. $v_p^2 = GM(2/r_p - 1/a)$ dengan $r_p =
8.77\times10^{10}$ m, $a = 2.67\times10^{12}$ m: $v_p = \qty{54}{km/s}$;
$v_a = v_pr_p/r_a = \qty{0.91}{km/s}$.
\end{solution}

\begin{solution}{exo:b1:central-forces:7}
$a = \qty{24490}{km}$. $v_1 = \qty{7.67}{km/s}$; pada elipsnya di $r_1$:
$\sqrt{GM(2/r_1 - 1/a)} = \qty{10.07}{km/s}$: $\Delta v_1 = \qty{2.4}{km/s}$.
Di $r_2$: $\sqrt{GM(2/r_2 - 1/a)} = \qty{1.62}{km/s}$, lingkarannya
\qty{3.07}{km/s}: $\Delta v_2 = \qty{1.45}{km/s}$. Totalnya \qty{3.85}{km/s};
durasinya $\pi\sqrt{a^3/GM} = \qty{1.9e4}{s} = \qty{5.3}{h}$.
\end{solution}

\begin{solution}{exo:b1:central-forces:8}
$E_{p,\mathrm{eff}}' = GMm/r^2 - L^2/mr^3 = 0$ pada $r_c = L^2/GMm^2$.
$E_{p,\mathrm{eff}}'' = -2GMm/r^3 + 3L^2/mr^4$; pada $r_c$, $L^2 = GMm^2r_c$
memberi $E'' = GMm/r_c^3$, sehingga $\omega_r^2 = E''/m = GM/r_c^3 = \omega_{\text{orb}}^2$.
Usikan radial yang kecil berosilasi tepat sekali tiap putaran: orbitnya
menutup pada dirinya sendiri --- sebuah elips yang sedikit lonjong, bukan roset.
\end{solution}

\begin{solution}{exo:b1:central-forces:9}
Hanya \omterm{def:b1:central-forces:gravitation}{gravitasi} yang bekerja pada stasiunnya dan awaknya sama rata;
keduanya jatuh pada orbit yang sama, sehingga tak diperlukan gaya sentuh
di antara keduanya --- tak ada lantai yang mendorong, dan karena itu
``tanpa berat''. $g = GM/r^2 = 3.99\times10^{14}/(6.77\times10^6)^2
= \qty{8.7}{m/s^2}$: \omterm{def:b1:central-forces:gravitation}{gravitasinya} nyaris sekuat penuh; beratnya tidak
lenyap, ia sepenuhnya terpakai untuk membengkokkan \omterm{def:b1:point-kinematics:frame}{lintasannya}.
\end{solution}

\begin{solution}{exo:b1:central-forces:10}
$a = 4\pi^2r/T^2 = 39.5 \times 3.84\times10^8/(2.36\times10^6)^2 =
\qty{2.7e-3}{m/s^2}$; $g(R/r)^2 = 9.81/60.3^2 = \qty{2.7e-3}{m/s^2}$: hukum
gaya yang sama yang menjatuhkan sebuah apel menahan Bulan, terencerkan sebagai $1/r^2$.
\end{solution}

\begin{solution}{exo:b1:central-forces:11}
$E_k = 2Ze^2/4\pi\varepsilon_0r_{\min}$: $r_{\min} = 8.99\times10^9 \times 2 \times 79
\times (1.60\times10^{-19})^2/8.0\times10^{-13} = \qty{4.5e-14}{m} = \qty{45}{fm}$,
enam kali jari-jari intinya: partikel $\alpha$-nya berbalik sebelum
menyentuh. Untuk mencapai \qty{7}{fm}: $E_k \approx \qty{32}{MeV}$.
\end{solution}

\begin{solution}{exo:b1:central-forces:12}
$E_m = -GMm/2r$ berkurang $\Rightarrow$ $r$ berkurang $\Rightarrow$ $v =
\sqrt{GM/r}$ bertambah. \omterm{thm:b1:work-and-energy:ke}{Energi kinetiknya} $GMm/2r$ naik sebesar $\abs{\Delta E}$
sementara \omterm{def:b1:work-and-energy:conservative}{energi potensialnya} $-GMm/r$ turun sebesar $2\abs{\Delta E}$:
separuh \omterm{def:b1:work-and-energy:conservative}{energi potensial} yang dilepaskan menjadi kinetik, dan separuhnya
dilesapkan sebagai kalor oleh \omterm{prop:b1:newton-dynamics:drag}{gaya hambatnya}.
\end{solution}

\begin{solution}{pb:b1:central-forces:1}
\textbf{1.} $T = \qty{43082}{s}$: $r = (GMT^2/4\pi^2)^{1/3} = (3.986\times10^{14}
\times 1.856\times10^9/39.48)^{1/3} = \qty{2.656e7}{m}$; ketinggiannya
\qty{20190}{km}.

\textbf{2.} $v = \sqrt{GM/r} = \qty{3.87}{km/s}$.

\textbf{3.} $a = v^2/r = GM/r^2 = \qty{0.565}{m/s^2} = g(R/r)^2$, tujuh
belas kali lebih lemah daripada di permukaannya --- tetapi tak nol:
satelitnya terus-menerus jatuh mengelilingi Bumi.

\textbf{4.} $E_m/m = -GM/2r = \qty{-7.5}{MJ/kg}$; $E_k/m = \qty{+7.5}{MJ/kg}$,
$E_p/m = \qty{-15.0}{MJ/kg}$.

\textbf{5.} Bumi berputar sekali tiap hari sideris relatif terhadap
bintangnya (kerangka orbitnya); setelah dua orbit satelitnya kembali di
atas titik tanah yang sama: jejak tanahnya berulang tiap hari sideris.

\textbf{6.} $\sin\beta = R/r = 0.240$: $\beta = 13.9^\circ$, yaitu cakram
selebar $28^\circ$.

\textbf{7.} Geostasionernya \qty{42200}{km}: $1.59$ kali lebih tinggi;
orbit navigasinya berada di bawahnya.

\textbf{8.} $GM(1/R - 1/2r) = 3.986\times10^{14}(1.570\times10^{-7} - 1.88\times
10^{-8}) = \qty{5.5e7}{J/kg} = \qty{55}{MJ/kg}$.

\textbf{9.} $r_1 = \qty{6771}{km}$: $v = \qty{7.67}{km/s}$.

\textbf{10.} $a = \qty{16665}{km}$; di $r_1$: $\sqrt{GM(2/r_1 - 1/a)} =
\qty{9.69}{km/s}$, $\Delta v_1 = \qty{2.0}{km/s}$; di $r_2$: \qty{2.47}{km/s}
terhadap lingkarannya \qty{3.87}{km/s}, $\Delta v_2 = \qty{1.4}{km/s}$; totalnya
\qty{3.4}{km/s}.

\textbf{11.} $\pi\sqrt{a^3/GM} = \qty{1.07e4}{s} \approx \qty{3.0}{h}$.

\textbf{12.} Lepas dari orbit rendah: $(\sqrt2 - 1)\times 7.67 = \qty{3.2}{km/s}$:
\qty{3.4}{km/s} misinya melampauinya --- mencapai orbit ini kira-kira
semahal meninggalkan Bumi sama sekali.

\textbf{13.} $m_0/m_1 = \eu^{3.4/3.0} = 3.1$: $m_1 = \qty{640}{kg}$,
propelannya \qty{1360}{kg} --- dua pertiga massanya.

\textbf{14.} $\cos\theta = 0.240$, $\theta = 76^\circ$; bagiannya $(1 - 0.240)/2
= 0.38$.

\textbf{15.} $24 \times 0.38 = 9$ rata-rata; empat terlampaui dengan
nyaman hampir selalu.

\textbf{16.} Zenit: $2.02\times10^7/c = \qty{67}{ms}$; cakrawala: jaraknya
$\sqrt{r^2 - R^2} = \qty{2.58e7}{m}$, \qty{86}{ms}.

\textbf{17.} $3/c = \qty{10}{ns}$.

\textbf{18.} Simpangan jam penerimanya menjadi besaran keempat yang tak
diketahui di samping ketiga koordinatnya: empat jarak, empat persamaan.

\textbf{19.} $10^{-8}\ \unit{s}/\qty{86400}{s} \approx \num{1e-13}$.

\textbf{20.} $v^2/2c^2 = (3874)^2/(2 \times 8.99\times10^{16}) = \num{8.3e-11}$:
$8.3\times10^{-11} \times 86400 = \qty{7.2}{\micro s}$ tiap hari, lambat.

\textbf{21.} $GM(1/R - 1/r)/c^2 = 3.986\times10^{14} \times 1.19\times10^{-7}/
8.99\times10^{16} = \num{5.3e-10}$: \qty{45.7}{\micro s} tiap hari, cepat.

\textbf{22.} Bersihnya $+\qty{38.5}{\micro s}$ tiap hari; $c \times 38.5\times10^{-6}
= \qty{11.5}{km}$ galat tiap hari.

\textbf{23.} $(10.23 - 10.22999999543)/10.23 = \num{4.47e-10}$, cocok
dengan nilai bersihnya $4.5\times10^{-10}$ (\qty{38.5}{\micro s} sepanjang \qty{86400}{s}).

\textbf{24.} Tonjolan khatulistiwa Bumi, tarikan Bulan dan Matahari,
tekanan radiasi surya (dan penyalaan pendorong yang kecil): segmen
tanahnya mengukur orbitnya terus-menerus lalu mengunggah efemeris yang
segar, yang disiarkan satelitnya.

\textbf{25.} Energi dan \omterm{def:b1:angular-momentum:L}{momentum sudut} menetapkan \omterm{prop:b1:central-forces:effective}{orbit lingkaran} dari
periodenya; Kepler III menempatkannya; geometri menetapkan cakupan dan
waktu cahayanya; dan jamnya --- jantung sistemnya --- harus dikoreksi
akibat kerelatifan sebesar \qty{38}{\micro s} sehari, atau penentuan
posisinya hanyut sejauh sepuluh kilometer menjelang malam.
\end{solution}
